\chapter{Messages on a grid: quadrature, truncation and their diagnostics}
\label{ch:grid-quadrature}

The recursion of \cref{ch:factor-graphs-bp} is exact for functions. A computer holds a
function by finitely many numbers, and the paper's choice is the simplest one: the
values of each message at $\ngrid$ equally spaced points of an interval
$[-\halfwidth,\halfwidth]$, with integrals replaced by the composite trapezoidal rule.
The paper summarises the consequence in one clause, ``exact for the discretised factor
model and controlled for the continuous one'', and devotes Appendix B to measuring the
control. This chapter explains the clause word by word: what the discretised model is
and why forward--backward is exact for it; where the trapezoidal rule's error comes
from, why it is far smaller than its textbook order for smooth rapidly decaying
integrands and why it is not for a discontinuous kernel; what happens when the domain
is too narrow for the data; and what each of the paper's diagnostics (column-mass
defect, boundary belief mass, drift, exceedance count) measures. It ends with the
normalisation bookkeeping of the implementation.

\section{Functions as vectors}

\subsection{The grid and its weights}

Fix a \term{half-width} $\halfwidth>0$ and a number of points $\ngrid\ge2$. The
\term{grid} is $\gridpt{k}=-\halfwidth+(k-1)\gridh$ for $k=1,\dots,\ngrid$ with
\term{spacing} $\gridh=2\halfwidth/(\ngrid-1)$; the paper's operating grid has
$\ngrid=401$, $\halfwidth=8$ and hence $\gridh=0.04$. A function $g$ on $\Reals$ is
represented by the vector $(g(\gridpt{1}),\dots,g(\gridpt{\ngrid}))$, and an integral
by the \term{composite trapezoidal rule}
\begin{equation}
  \int_{-\halfwidth}^{\halfwidth}g(a)\dd a\;\approx\;\sum_{k=1}^{\ngrid}w_k\,g(\gridpt{k}),
  \qquad w_1=w_{\ngrid}=\tfrac{\gridh}{2},\quad w_k=\gridh\ \text{otherwise}.
  \label{eq:gq-trap}
\end{equation}
The weights come from integrating the piecewise linear interpolant of $g$ through the
grid values: each interior cell $[\gridpt{k},\gridpt{k+1}]$ contributes
$\gridh\,(g(\gridpt{k})+g(\gridpt{k+1}))/2$, and summing over cells gives every interior
point the weight $\gridh$ and the two end points half of it. The code is
\code{make\_grid} in \code{src/\allowbreak markov\_score/bp\_grid.py}.

\subsection{The discretised factor model}

Replacing the integrals of the chain posterior by \cref{eq:gq-trap} amounts to
replacing the continuous model by a finite-state one. Define the
\term{discretised factor model} as the probability mass function on
$\{1,\dots,\ngrid\}^{\nsites}$
\begin{equation}
  \hat p(k_1,\dots,k_\nsites\mid\vx)\propto
  \mu(\gridpt{k_1})w_{k_1}\prod_{i=2}^{\nsites}\kernel(\gridpt{k_i}\mid\gridpt{k_{i-1}})w_{k_i}
  \prod_{i=1}^{\nsites}\like(x_i\mid\gridpt{k_i}).
  \label{eq:gq-discretised}
\end{equation}
It is a hidden Markov model with $\ngrid$ states whose transition matrix is
$\hat\kernel[k,l]=\kernel(\gridpt{k}\mid\gridpt{l})\,w_k$, that is, the continuous
kernel sampled on the grid and multiplied by the quadrature weight of the destination
state. Two remarks. First, $\hat\kernel$ is not exactly stochastic: its column sums
$\sum_k\hat\kernel[k,l]$ are trapezoidal estimates of $\int\kernel(a\mid\gridpt{l})\dd a=1$
and miss one by the quadrature and truncation error of that integral. The recursion
normalises every message, so this deficit never accumulates, but it is the natural
scalar diagnostic of the discretisation at its source and is measured as the
\term{column-mass defect} $\max_l\bigl|\sum_k\kernel(\gridpt{k}\mid\gridpt{l})w_k-1\bigr|$
(\cref{sec:gq-diagnostics}). Second, running the two sweeps of \cref{ch:factor-graphs-bp}
with the matrix $\hat\kernel$ is the forward--backward algorithm on this finite-state
model, and by the exactness theorem of that chapter it returns the marginals and the
normaliser of \cref{eq:gq-discretised} \emph{exactly}, up to floating point. The
returned evidence is $\log\hat Z(\vx)$, the normaliser of the discretised model: exact
bookkeeping of an approximate object. This is the first half of the paper's clause.
The tests \code{test\_xi\_matches\_brute\_force\_enumeration} in
\code{tests/test\_em\_bp.py} and \code{tests/\allowbreak test\_fisher\_bruteforce.py} check exactly
this statement, by enumerating \cref{eq:gq-discretised} on a tiny grid.

The second half, ``controlled for the continuous one'', concerns the distance between
the marginals of \cref{eq:gq-discretised}, read as functions on the grid, and those of
the continuous posterior. That distance has two sources, treated in turn:
\term{quadrature error}, from replacing integrals over $[-\halfwidth,\halfwidth]$ by
sums, and \term{truncation error}, from ignoring the part of $\Reals$ outside the
domain.

\section{The trapezoidal rule and its error}

\subsection{Euler--Maclaurin}

Let $g$ be $2m$ times continuously differentiable on $[-\halfwidth,\halfwidth]$ and
write $T_\gridh(g)$ for the right-hand side of \cref{eq:gq-trap}. The
\term{Euler--Maclaurin formula} states
\begin{equation}
\begin{aligned}
  T_\gridh(g)-\int_{-\halfwidth}^{\halfwidth}g\dd a
  &=\sum_{j=1}^{m-1}\frac{B_{2j}}{(2j)!}\gridh^{2j}\Bigl[g^{(2j-1)}(\halfwidth)-g^{(2j-1)}(-\halfwidth)\Bigr]+R_m,\\
  |R_m|&\le C_m\,\gridh^{2m}\int_{-\halfwidth}^{\halfwidth}|g^{(2m)}|\dd a,
\end{aligned}
  \label{eq:gq-em}
\end{equation}
where $B_2=\tfrac16$, $B_4=-\tfrac1{30}$, $B_6=\tfrac1{42}$, \dots\ are the Bernoulli
numbers and $C_m$ is an absolute constant. The derivation for one cell is worth
seeing once. On $[0,\gridh]$, integrate by parts with the polynomial
$P_1(a)=a-\gridh/2$ (whose derivative is $1$ and which is odd about the midpoint):
$\int_0^\gridh g\dd a=\bigl[P_1 g\bigr]_0^\gridh-\int_0^\gridh P_1g'\dd a
=\tfrac\gridh2\bigl(g(0)+g(\gridh)\bigr)-\int_0^\gridh P_1g'\dd a$. The first term is the
trapezoidal rule on the cell; the second is the error. Integrating by parts again with
$P_2$ chosen so that $P_2'=P_1$ and $P_2(0)=P_2(\gridh)=\gridh^2/12$ turns it into
$-\tfrac{\gridh^2}{12}\bigl(g'(\gridh)-g'(0)\bigr)+\int_0^\gridh P_2g''\dd a$, and
summing over the cells telescopes the boundary terms to the two ends of
$[-\halfwidth,\halfwidth]$, giving the $j=1$ term of \cref{eq:gq-em} with
$B_2/2!=1/12$. Continuing the integrations by parts generates the higher Bernoulli
terms. The textbook conclusion is that the trapezoidal rule is \term{second order}:
its error is $O(\gridh^2)$, with constant $\tfrac1{12}|g'(\halfwidth)-g'(-\halfwidth)|$.
The paper's Appendix B calls the quadrature error ``second order in the spacing'' for
this reason.

\subsection{Why the rule is much better than second order on a Gaussian}

Look again at \cref{eq:gq-em}. Every term of the sum involves derivatives of $g$ at the
end points of the domain. If $g$ and all its derivatives are negligible there, as they
are for a Gaussian-like integrand whose domain is several standard deviations wide,
every power of $\gridh$ in the expansion has a vanishing coefficient, and the error is
smaller than any power of $\gridh$. This property is called \term{spectral accuracy}
(or super-algebraic convergence). It can be made exact for an integrand on the whole
line: for $g$ smooth and rapidly decaying, the \term{Poisson summation formula} gives
\begin{equation}
  \gridh\sum_{k\in\mathbb{Z}}g(k\gridh)-\int_\Reals g\dd a
  =\sum_{n\ne0}\hat g\!\left(\frac{2\pi n}{\gridh}\right),\qquad
  \hat g(\omega)=\int_\Reals g(a)e^{-i\omega a}\dd a,
  \label{eq:gq-poisson}
\end{equation}
so the error of the infinite trapezoidal sum is the sum of the Fourier transform of
$g$ at the frequencies $2\pi n/\gridh$: it decays as fast as $\hat g$ does. For the
standard Gaussian density, $\hat g(\omega)=e^{-\omega^2/2}$ and the error is
$2\sum_{n\ge1}e^{-2\pi^2n^2/\gridh^2}$, dominated by $2e^{-2\pi^2/\gridh^2}$. At
$\gridh=1$ this is $2e^{-19.74}\approx5\times10^{-9}$; at $\gridh=0.5$ it is
$2e^{-79}\approx10^{-34}$, below double precision; at the operating $\gridh=0.04$ it
is zero for every practical purpose. The worked example of \cref{sec:gq-example}
checks the $\gridh=1$ figure on a finite domain.

This is why the quadrature is not the binding error in the paper's numerics for smooth
families, and why the paper writes that ``the domain, not the spacing, is what a
heavy-tailed family tests''. Two situations degrade the picture. If the integrand is
narrow compared with $\gridh$, which happens to the unary factor
$\like(x_i\mid\cdot)$ at small $t$ (its standard deviation as a function of $a$ is
$\sqrt{\Dt}/e^{-t}$, about $0.32$ at $t=0.05$, eight cells at $\gridh=0.04$), the
Fourier transform decays slowly and the terms of \cref{eq:gq-poisson} are no longer
negligible; the code exposes \code{likelihood\_resolution\_ok} in
\code{src/\allowbreak markov\_score/noising.py} for this check, and the paper's E04 refines the
grid on the reported cells to show the effect is small. And if $g$ is not smooth the
expansion \cref{eq:gq-em} fails at the first derivative that does not exist: the
Laplace kernel has a kink at $a'=\corr a$ (its first derivative jumps), which costs
the Bernoulli terms beyond $j=1$ but still leaves an $O(\gridh^2)$ error; the uniform
kernel is discontinuous, and that is the case of the next subsection.

\subsection{A discontinuous kernel is first order}
\label{sec:gq-discontinuous}

The uniform innovation has kernel $\kernel(a'\mid a)=1/(2h_u)$ for
$|a'-\corr a|\le h_u$ and zero otherwise, with $h_u=\sqrt{3\ivar}$ so that the variance
is $\ivar$. As a function of the integration variable it is a step, and the trapezoidal
rule on a step is only \term{first order}. The mechanism is elementary. If a jump of
height $J$ falls at a distance $\delta\in[0,\gridh)$ from the nearest grid point, the
piecewise linear interpolant misplaces the jump by up to $\gridh$ and the integral by
an amount between $-J\gridh/2$ and $+J\gridh/2$, depending on $\delta$; the error
oscillates with the position of the jump and its envelope shrinks only like $\gridh$,
not $\gridh^2$. \Cref{sec:gq-example} exhibits the error exactly in the aligned case.
In the chain the jumps sit at $a'=\corr a\pm h_u$ and move with $a$, so every column of
$\hat\kernel$ has its own misplacement, and the paper's Appendix B records the
signature: for the uniform family the worst drift is largest on the \emph{coarsest}
spacing rather than on the narrowest domain, the opposite of the heavy-tailed families.
The column-mass defect of the uniform kernel is likewise the largest among the
families at moderate correlation, and it does fall with $\ngrid$.

\section{Truncation}

\subsection{A site beyond the domain}

The second error is the one the paper's sweep found to be binding. The grid represents
only $[-\halfwidth,\halfwidth]$, so every message is implicitly zero outside. For a
family whose marginal has light tails (Gaussian, Laplace, the bimodal mixture, the
generalised Gaussian) the probability that any site of any sequence lies beyond
$\halfwidth=8$ is negligible: for a standard Gaussian site it is about $10^{-15}$, and
for the Laplace innovation at the frozen correlation the probability that a single
innovation exceeds $8$ in absolute value is of order $10^{-10}$. For a Student-$t$
innovation the tail is polynomial and the probability is not negligible, which is the
case treated next.

Suppose a clean site has $|a_i|>\halfwidth$. Its observation is
$x_i=e^{-t}a_i+\sqrt{\Dt}z_i$, and the unary factor $\like(x_i\mid a)$, as a function of
the latent $a$, is a Gaussian centred at $x_i/e^{-t}\approx a_i$, outside the grid.
Restricted to the grid the factor is a monotone ramp, largest at the boundary point
nearest to $a_i$; the transition factors from the neighbouring sites add little
against it. The belief at site $i$ is therefore concentrated in the last few cells at
the edge: the paper's phrase is that the posterior is \term{pinned to the boundary}.
The posterior mean is then close to $\pm\halfwidth$ instead of close to $a_i$, and by
Tweedie's identity the score error at that site is about
$(e^{-t}/\Dt)\,(|a_i|-\halfwidth)$, which at the lowest level of the schedule is a
factor $e^{-t}/\Dt\approx10$ times the excess. One such site in a bundle of thousands
can dominate a pooled relative error, and, crucially, it is pinned on \emph{every}
grid with the same half-width, so refining $\ngrid$ at fixed $\halfwidth$ does not
move it. That is the signature the paper's Appendix B reports for the binding cell:
drift flat in $\ngrid$, falling only when the domain widens.

\subsection{The exceedance count and its Poisson law}

The right summary of truncation for a polynomial tail is not a maximum over sites but
a count. Let $N_{\halfwidth}$ be the number of sites among the $\nseq\nsites$ of a
bundle with $|a_i|>\halfwidth$. Its expectation is
\begin{equation}
  \E N_{\halfwidth}=\nseq\nsites\,\Prob(|a|>\halfwidth)
  \label{eq:gq-exceed}
\end{equation}
when every site has the same marginal law (at $\corr=0$ every site after the first is
an independent draw of the innovation). Because the events are rare and, along a
chain, only weakly dependent, the count is close to a Poisson variable with that mean:
this is the law of rare events (Poisson's limit theorem, or Le Cam's inequality for
dependent indicators with small total variation). No finite $\halfwidth$ makes
\cref{eq:gq-exceed} zero for a polynomial tail; widening the domain only lowers the
mean, and a Poisson variable with mean $\lambda$ is zero with probability $e^{-\lambda}$
and at least one with probability $1-e^{-\lambda}$. For the Student-$t$ innovation with
$\nu$ degrees of freedom, variance-matched so that $a=s\,T_\nu$ with
$s^2=\ivar(\nu-2)/\nu$, the tail probability is
$\Prob(|a|>\halfwidth)=2\,\bar F_\nu(\halfwidth/s)$ with $\bar F_\nu$ the survival
function of $T_\nu$; it decays like $\halfwidth^{-\nu}$. \Cref{sec:gq-example}
computes the mean count for the paper's grids. The paper's Appendix B states this law
and gives the expected counts at $\halfwidth=12$ and $16$; the point is that the
Student-$t$ cell of the family ladder is affected by a handful of sites, whose number
is predicted, and not by ``the family'' in some diffuse sense.

\section{The diagnostics}
\label{sec:gq-diagnostics}

The paper measures four quantities, all produced by \code{experiments/\allowbreak e11\_grid\_validation.py}
and the diagnostic branch of \code{grid\_bp\_batch}. For each the definition,
what it detects, and what it does not.

\paragraph{Column-mass defect.}
$\max_l\bigl|\sum_k\kernel(\gridpt{k}\mid\gridpt{l})w_k-1\bigr|$: the worst deviation of a
column of $\hat\kernel$ from unit mass. It combines two effects. In the interior of the
grid it is the quadrature error of a one-dimensional integral of the kernel, spectrally
small for a smooth kernel (machine precision for the Gaussian and mixture families at
the operating grid) and $O(\gridh)$ for the uniform one. In the columns near the
boundary it is a \emph{truncation} effect: the column with parent $\gridpt{l}=\pm\halfwidth$
is a density centred at $\pm\corr\halfwidth$, and the part of it beyond $\pm\halfwidth$
is cut off. As $\corr$ rises, $\corr\halfwidth$ moves towards the edge and the cut-off
mass grows, from nothing at $\corr=0$ to a sizeable fraction of a column at
$\corr=0.95$; this part is flat in $\ngrid$ and falls only with $\halfwidth$. A reader
checking the paper's Appendix B should therefore read the growth of the defect with
correlation as a statement about the outermost columns, not about resolution in the
interior, where the smooth families are resolved to machine precision at every
correlation. In either case the defect is not an error bar for the posterior: the
recursion renormalises every message, so a defect that is the same across a column's
neighbours cancels in the normalised beliefs, and the paper reports it as a diagnostic
only.

\paragraph{Boundary belief mass.}
For each site of each chain, the mass of the normalised belief (or of the forward
message) in the outermost $2\%$ of the grid on each side,
$\sum_{k\le e}b_i(\gridpt{k})w_k+\sum_{k>\ngrid-e}b_i(\gridpt{k})w_k$ with
$e=\lfloor0.02\,\ngrid\rfloor$, maximised over sites and chains. It is computed inside the
recursion, on the very objects the recursion propagates, which is the reason the
diagnostic lives in \code{grid\_bp\_batch} rather than in a separate tool: an earlier
version computed a ``boundary mass'' from the likelihood table alone, which involved
no message and measured nothing the algorithm used. A value near one means that some
site's posterior is pinned to the boundary, and the count of such sites is the
exceedance count above. A value of a few per cent at $\corr=0.85$ with no site beyond
the domain is a different, benign situation: a site near $|a|\approx6$ whose posterior
at the lowest level has a tail reaching the edge. The paper quotes the frozen
configuration's value and the binding cell's value in Appendix B.

\paragraph{Drift.}
The relative root-mean-square difference between the posterior-mean field computed on
a candidate grid and on the finest grid of the ladder, on the same chains and noise
draws, pooled over the twelve levels (and separately at the worst level and at the
lowest level). Drift is what a reported number would move by under refinement, and it
is the quantity a claim is judged against; the other diagnostics say why it is small.
Its reference is the finest grid, not the truth, so it is a self-convergence
diagnostic: a site pinned on both grids contributes no drift, which is why the
binding cell's reference grid in E11 could itself miss the largest site and why E18
and the wide-domain reruns E07b and E07c widen the reference further.

\paragraph{Exceedance count.}
Not computed by the sweep but by \cref{eq:gq-exceed} from the family's tail; it is the
interpretation of the previous three for heavy tails, and it was the audit's re-reading
of the E11 cell that turned ``all the belief mass outside the domain'' into ``one or two
sites of eight thousand''.

\subsection{How the experiments use them}

\emph{E11 (the envelope).} Five families, four correlations, twelve levels, four seeds,
seven candidate grids against a reference at a finer and wider grid. Its findings are
the paper's Appendix B: at the frozen configuration the drift is far below every
reported effect; across the sweep the worst drift is flat in $\ngrid$ at fixed
$\halfwidth$ and falls only when $\halfwidth$ grows (truncation, in the Student-$t$
cell at $\corr=0$); the column-mass defect grows with correlation but does not
predict the drift; the uniform family is the one case whose worst drift sits on the
coarsest spacing. \emph{E11b} repeats the sweep at chain lengths $64$ and $128$ to
show that the envelope measured at $\nsites=32$ covers the longer chains of the length
study; its prediction is that per-site errors compound at most linearly in $\nsites$.
\emph{E07b and E07c} recompute the family ladder of the paper's Appendix F on wider
domains ($\halfwidth=12$ and $16$), which by \cref{eq:gq-exceed} lowers the expected
exceedance count by an order of magnitude and then by another. \emph{E18} recomputes
the closure-gap panel of the paper's Figure 2 at $(\ngrid,\halfwidth)=(1601,16)$ and
$(3201,32)$ and shows that the Student-$t$ departure from the $e^{-8t}$ law on the
operating grid is the truncated prior's tail in the reference, not the moment
condition; on the wide domain the family reaches the predicted constant. \emph{E04}
re-scores the validation-selected fitted kernels of the capacity experiment at
$\ngrid=801$, $1601$ and at $\halfwidth=10$, with the reference recomputed at each
resolution, to bound the movement of the reported estimands under refinement; the
paper's Appendix B quotes the bound and notes it is orders of magnitude below the
effects.

\section{Normalisation and the evidence on the grid}
\label{sec:gq-normalisation}

The implementation in \code{src/\allowbreak markov\_score/bp\_grid.py} keeps messages in the
linear domain and rescales them after every step, accumulating the logarithms of the
scale factors. The details matter for reading the code and for trusting the returned
log-evidence.

\paragraph{The likelihood table.}
The unary factors are tabulated with the row maximum removed,
\[
  \log\like(x_i\mid\gridpt{k})\ \text{is stored as}\ -\frac{(x_i-e^{-t}\gridpt{k})^2}{2\Dt}+c_i,\qquad
  c_i=\max_k\frac{(x_i-e^{-t}\gridpt{k})^2}{2\Dt},
\]
so that the largest entry of every row is exactly $1$ after exponentiation and no row can
underflow to zero (\code{log\_likelihood\_matrix} in \code{src/\allowbreak markov\_score/noising.py};
\code{grid\_bp\_batch} recomputes the same table internally). The Gaussian constant
$-\tfrac12\log(2\pi\Dt)$ is dropped at the same time. Both removals are harmless for
the beliefs, which are normalised, and both are recorded: the sum of the row shifts
and $\nsites$ copies of the constant are added back to the evidence at the end.

\paragraph{The sweeps.}
With $\hat\kernel$ the matrix $\kernel(\gridpt{k}\mid\gridpt{l})$ (weights applied on the
incoming side) the forward step is
$\lmsg_{i+1}=\hat\kernel\,(\lmsg_i\circ\ell_i\circ w)/c_{i+1}$, where
$c_{i+1}=\sum_k w_k[\hat\kernel(\lmsg_i\circ\ell_i\circ w)]_k$ is the trapezoidal mass of
the unnormalised result; the backward step is the same with $\hat\kernel^{\!\top}$ and
its own constants, which are not needed for the evidence. Each step is a
matrix--vector product, or, for a batch of $B$ chains, a matrix--matrix product; the
batch is processed in chunks of at most $256$ chains so that the
$(\nsites\times B\times\ngrid)$ temporaries fit in memory, with results identical to
round-off across chunk sizes (\code{tests/\allowbreak test\_bp\_grid\_chunk.py}). A guard raises
an error if any normaliser is non-finite or non-positive: a message that has lost all
its mass is a failure to report, not to hide.

\paragraph{The evidence.}
Putting the pieces together,
\begin{equation}
  \log\hat Z(\vx)=\sum_{i=2}^{\nsites}\log c_i
  +\log\sum_k w_k\,\lmsg_\nsites(\gridpt{k})\,\ell_\nsites(\gridpt{k})
  +\sum_{i=1}^{\nsites}\bigl(-c_i^{\mathrm{shift}}\bigr)-\frac{\nsites}{2}\log(2\pi\Dt),
  \label{eq:gq-evidence}
\end{equation}
which is the log-normaliser of \cref{eq:gq-discretised} to floating point. Its
correctness is checked in three independent ways: against the closed-form Gaussian
marginal likelihood; against an enumeration of the discretised posterior; and against
a Monte Carlo estimate on the Laplace chain, where no closed form exists. The first
and the third are tests in \code{tests/test\_em\_bp.py}; the second is the file
\code{tests/test\_fisher\_bruteforce.py}. Its role in the paper is twofold: it is the
objective that expectation--maximisation increases (\cref{ch:em-through-noise}), and,
summed over a held-out one-view bundle and divided by the number of transitions, it
is the oracle-free selection rule of E20.

\paragraph{Diagnostic outputs.}
When called with \code{boundary\_frac}, \code{grid\_bp\_batch} also returns a dictionary
with the maximum and mean boundary mass of the forward messages and of the beliefs,
and the number of edge cells used. These are the numbers behind the paper's Appendix B
and behind the frozen configuration's docstring in \code{experiments/frozen\_config.py},
which is careful to say that a maximum over sampled sites is a statistic and not a
bound.

\section{Worked examples}
\label{sec:gq-example}

\paragraph{The Gaussian, spectral accuracy.}
Integrate the standard Gaussian density on $[-\halfwidth,\halfwidth]$ with the
trapezoidal rule. At $\halfwidth=8$ and $\ngrid=17$, so $\gridh=1$, the nine
non-negligible grid values are $\phi(0)=0.398942$, $\phi(\pm1)=0.241971$,
$\phi(\pm2)=0.053991$, $\phi(\pm3)=0.004432$, $\phi(\pm4)=0.000134$, and the sum
$\gridh\sum_k\phi(k)$ (the end weights are irrelevant at this precision) is
$1.0000000054$: an error of $5\times10^{-9}$, which is $2e^{-2\pi^2}$ to two digits,
as \cref{eq:gq-poisson} predicts. Halving the spacing to $\gridh=0.5$ gives an error
below $10^{-15}$. By contrast, at $\gridh=2$ the sum is $1.0144$: the Poisson term
$2e^{-2\pi^2/4}=0.0144$ again. The convergence is not $\gridh^2$; it is exponential in
$1/\gridh^2$.

\paragraph{The step, first order.}
Integrate the uniform density $\tfrac12\mathbf 1\{|a|\le1\}$ on $[-2,2]$ with
$\ngrid=9$ points, $\gridh=0.5$. The grid values at $a=\pm1$ are $\tfrac12$ (the jumps
sit on grid points), at $|a|<1$ they are $\tfrac12$, elsewhere zero; the trapezoidal
sum is $\gridh\cdot\tfrac12\cdot5=1.25$, an error of $0.25=\gridh/2$. At $\gridh=0.25$
the error is $0.125=\gridh/2$; at $\gridh=0.125$ it is $0.0625$. The error halves when
$\gridh$ halves: first order, with constant $J/2$ per jump, exactly as
\cref{sec:gq-discontinuous} says. Shifting the grid so that the jumps fall between
grid points changes the sign and the size of the error but not its order.

\paragraph{The column-mass defect on a coarse grid.}
Take the Gaussian kernel $\kernel(a'\mid a)=\normal{\corr a}{\ivar}$ with $\corr=0.5$,
$\ivar=0.75$ on the five-point grid $\{-2,-1,0,1,2\}$, $\gridh=1$. The column with parent
$\gridpt{l}=0$ is the density $\normal{0}{0.75}$ sampled at the five points,
$(0.0320,\,0.2365,\,0.4607,\,0.2365,\,0.0320)$, and its trapezoidal mass with end weights
$\tfrac12$ is $0.9657$: a defect of $0.034$, almost all of it the tail of a
$\normal{0}{0.75}$ beyond $\pm2$, which is $0.021$, plus the quadrature error of
$\gridh=1$ on a Gaussian of standard deviation $0.87$. The column with parent $\gridpt{l}=2$
is centred at $1$ and loses the whole of its tail beyond $2$; its mass is $0.848$, a
defect of $0.152$ that no refinement of $\gridh$ removes and that a wider domain does.
Widening to $[-4,4]$ at the same spacing brings the worst column to $0.983$. This is
the mechanism behind the growth of the paper's defect with correlation: the extreme
column's centre $\corr\halfwidth$ approaches the edge.

\paragraph{The exceedance count.}
For the variance-matched Student-$t_5$ innovation, $s^2=0.6$ and
$\Prob(|a|>\halfwidth)=2\bar F_5(\halfwidth/\sqrt{0.6})$. At $\halfwidth=8$ the argument is
$10.33$ and the two-sided tail is $1.46\times10^{-4}$; over $\nseq\nsites=256\times32=8192$
sites the expected count is $1.2$, so the probability of no exceedance in a bundle is
$e^{-1.2}\approx0.30$ and of two or more is $0.34$. At $\halfwidth=12$ the expected
count is $0.17$ and at $\halfwidth=16$ it is $0.04$. These are the analytic figures the
paper's Appendix B quotes next to the observed counts of the E11 cell; a reader can
reproduce them from the $t$ survival function alone.

\where{The paper, Section 2 (``messages are functions on $\Reals$; we carry them on
a grid \dots exact for the discretised factor model and controlled for the continuous
one'') and Appendix B (the three levels of exactness, the validation at the frozen
configuration, the envelope, the refinement, the floors); Appendix F (the Student-$t$
cell excluded pending a wider domain); Appendix E (E18, the wide-domain closure
panel). Experiments E11, E11b, E07b, E07c, E18 and E04. Code: the grid and the two
sweeps live in \code{src/\allowbreak markov\_score/bp\_grid.py} (\code{make\_grid}, \code{grid\_bp},
and \code{grid\_bp\_batch} with its \code{chunk} and \code{boundary\_frac} arguments);
the row-shifted likelihood table and \code{likelihood\_resolution\_ok} in
\code{src/\allowbreak markov\_score/noising.py}; the drivers are
\code{experiments/\allowbreak e11\_grid\_validation.py}, \code{experiments/\allowbreak e18\_closure\_wide.py},
\code{experiments/\allowbreak e04\_refinement.py} and \code{experiments/\allowbreak e07\_phase\_diagram.py},
with the macros emitted by \code{tools/\allowbreak make\_grid\_envelope.py}. Tests: for the grid
itself \code{tests/\allowbreak test\_grid\_convergence.py} and \code{tests/\allowbreak test\_bp\_grid\_chunk.py};
for the evidence, \code{tests/test\_em\_bp.py} and \code{tests/\allowbreak test\_fisher\_bruteforce.py}.}

\paragraph{References.}
P.~J.~Davis and P.~Rabinowitz (1984), \emph{Methods of Numerical Integration}, 2nd ed.,
Academic Press, Chapters 2 and 4.
L.~N.~Trefethen and J.~A.~C.~Weideman (2014), ``The exponentially convergent trapezoidal
rule'', \emph{SIAM Review} 56(3).
G.~Kitagawa (1987), ``Non-Gaussian state-space modeling of nonstationary time series'',
\emph{Journal of the American Statistical Association} 82(400).
L.~Le~Cam (1960), ``An approximation theorem for the Poisson binomial distribution'',
\emph{Pacific Journal of Mathematics} 10(4).
N.~J.~Higham (2002), \emph{Accuracy and Stability of Numerical Algorithms}, 2nd ed., SIAM,
Chapter 2 (floating-point arithmetic and rescaling).
